\documentclass[10pt,letterpaper]{article}
\usepackage[utf8]{inputenc}
\usepackage[intlimits]{amsmath}
\usepackage{amsfonts}
\usepackage{amssymb}
\usepackage{ragged2e}
\usepackage[letterpaper, margin=0.5in]{geometry}
\usepackage{graphicx}
\usepackage{cancel}
\usepackage{mathtools}
\usepackage{tabularx}
\usepackage{arydshln}
\usepackage{tensor}
\usepackage{array}
\usepackage{xcolor}
\usepackage[boxed]{algorithm}
\usepackage[noend]{algpseudocode}
\usepackage{listings}
\usepackage{textcomp}
% \usepackage[pdf,tmpdir,singlefile]{graphviz}
\usepackage{mathrsfs}
\usepackage{bbm}
\usepackage{tikz}
\usepackage{tikz-cd}
\usepackage{enumitem}
\usepackage{arydshln}
\usepackage{relsize}
\usepackage{multicol}
\usepackage{scalerel}
\usepackage{upgreek}
\usepackage{wasysym}
\usepackage{ifthen}

\usetikzlibrary{bayesnet}
\setlist{noitemsep}

%%%%%%%%%%%%%%%%%%%%%%%%%%%%%
% Formatting commands
%%%%%%%%%%%%%%%%%%%%%%%%%%%%%
\newcommand{\n}{\hfill\break}
\newcommand{\nn}{\vspace{0.5\baselineskip}\n}
\newcommand{\up}{\vspace{-\baselineskip}}
\newcommand{\hangblock}[2]{\par\noindent\settowidth{\hangindent}{\textbf{#1: }}\textbf{#1: }\nolinebreak#2}
\newcommand{\lemma}[1]{\hangblock{Lemma}{#1}}
\newcommand{\defn}[1]{\hangblock{Defn}{#1}}
\newcommand{\thm}[1]{\hangblock{Thm}{#1}}
\newcommand{\cor}[1]{\hangblock{Cor}{#1}}
\newcommand{\prop}[1]{\hangblock{Prop}{#1}}
\newcommand{\ex}[1]{\hangblock{Ex}{#1}}
\newcommand{\exer}[1]{\hangblock{Exer}{#1}}
\newcommand{\fact}[1]{\hangblock{Fact}{#1}}
\newcommand{\remark}[1]{\hangblock{Remark}{#1}}
\newcommand{\proven}{\;$\square$\n}
\newcommand{\problem}[1]{\par\noindent{\nolinebreak#1}\n}
\newcommand{\problempart}[2]{\par\noindent\indent{}\settowidth{\hangindent}{\textbf{(#1)} \indent{}}\textbf{(#1) }\nolinebreak#2\n}
\newcommand{\ptxt}[1]{\textrm{\textnormal{#1}}}
\newcommand{\inlineeq}[1]{\centerline{$\displaystyle #1$}}
\newcommand{\pageline}{\up\par\noindent\rule{\textwidth}{0.1pt}}

%%%%%%%%%%%%%%%%%%%%%%%%%%%%%
% Math commands
%%%%%%%%%%%%%%%%%%%%%%%%%%%%%
% Set Theory
\newcommand{\card}[1]{\left|#1\right|}
\newcommand{\set}[1]{\left\{#1\right\}}
\newcommand{\setmid}{\;\middle|\;}
\newcommand{\ps}[1]{\mathcal{P}\left(#1\right)}
\newcommand{\pfinite}[1]{\mathcal{P}^{\ptxt{finite}}\left(#1\right)}
\newcommand{\naturals}{\mathbb{N}}
\newcommand{\N}{\naturals}
\newcommand{\integers}{\mathbb{Z}}
\newcommand{\Z}{\integers}
\newcommand{\rationals}{\mathbb{Q}}
\newcommand{\Q}{\rationals}
\newcommand{\reals}{\mathbb{R}}
\newcommand{\R}{\reals}
\newcommand{\complex}{\mathbb{C}}
\newcommand{\C}{\complex}
\newcommand{\halfPlane}{\mathbb{H}}
\let\HSym\H
\let\H\relax
\newcommand{\H}{\halfPlane}
\newcommand{\comp}{^{\complement}}
\DeclareMathOperator{\Hom}{Hom}
\newcommand{\Ind}{\mathbbm{1}}
\newcommand{\cut}{\setminus}
\DeclareMathOperator{\elem}{elem}

% Differentiable Manifolds
\newcommand{\RP}{\mathbb{RP}}
\newcommand{\CP}{\mathbb{RP}}
\newcommand{\osubset}{\overset{\mathclap{\scalebox{0.5}{\ptxt{open}}}}{\subset}}
\newcommand{\osubseteq}{\overset{\mathclap{\scalebox{0.5}{\ptxt{open}}}}{\subseteq}}
\newcommand{\osupset}{\overset{\mathclap{\scalebox{0.5}{\ptxt{open}}}}{\supset}}
\newcommand{\osupseteq}{\overset{\mathclap{\scalebox{0.5}{\ptxt{open}}}}{\supseteq}}
\newcommand{\pdat}[3]{\left.\pd{#1}{#2}\right|_{#3}}
\DeclareMathOperator{\so}{so}
\DeclareMathOperator{\codim}{codim}
\DeclareMathOperator{\Diff}{Diff}
\let\dSym\d
\let\d\relax
\newcommand{\d}{\partial}
\DeclareMathOperator{\gl}{gl}

% Graph Theory
\let\deg\relax
\DeclareMathOperator{\deg}{deg}
\newcommand{\degp}{\ptxt{deg}^{+}}
\newcommand{\degn}{\ptxt{deg}^{-}}
\newcommand{\precdot}{\mathrel{\ooalign{$\prec$\cr\hidewidth\hbox{$\cdot\mkern0.5mu$}\cr}}}
\newcommand{\succdot}{\mathrel{\ooalign{$\cdot\mkern0.5mu$\cr\hidewidth\hbox{$\succ$}\cr\phantom{$\succ$}}}}
\DeclareMathOperator{\cl}{cl}
\DeclareMathOperator{\affdim}{affdim}

% Probability
\newcommand{\parSymbol}{\P}
\newcommand{\Prob}{\mathbb{P}}
\renewcommand{\P}{\Prob}
\newcommand{\Avg}{\mathbb{E}}
\newcommand{\E}{\Avg}
\DeclareMathOperator{\Var}{Var}
\DeclareMathOperator{\cov}{cov}
\DeclareMathOperator{\Unif}{Unif}
\DeclareMathOperator{\Binom}{Binom}
\newcommand{\CI}{\mathrel{\text{\scalebox{1.07}{$\perp\mkern-10mu\perp$}}}}

% Standard Math
\newcommand{\inv}{^{-1}}
\newcommand{\abs}[1]{\left|#1\right|}
\newcommand{\ceil}[1]{\left\lceil{}#1\right\rceil{}}
\newcommand{\floor}[1]{\left\lfloor{}#1\right\rfloor{}}
\newcommand{\conj}[1]{\overline{#1}}
\newcommand{\of}{\circ}
\newcommand{\tri}{\triangle}
\newcommand{\inj}{\hookrightarrow}
\newcommand{\surj}{\twoheadrightarrow}
\newcommand{\ndiv}{\nmid}
\renewcommand{\epsilon}{\varepsilon}
\newcommand{\divides}{\mid}
\newcommand{\ndivides}{\nmid}
\DeclareMathOperator{\lcm}{lcm}
\DeclareMathOperator{\sgn}{sgn}
\newcommand{\map}[4]{\!\!\!\begin{array}[t]{rcl}#1 & \!\!\!\!\to & \!\!\!\!#2\\ {}#3 & \!\!\!\!\mapsto & \!\!\!\!#4\end{array}}
\newcommand{\bigsum}[2]{\smashoperator[lr]{\sum_{\scalebox{#1}{$#2$}}}}
\DeclareMathOperator{\gcf}{gcf}
\newcommand{\restr}[2]{\left.#1\right|_{#2}}

% Linear Algebra
\newcommand{\Id}{\textrm{\textnormal{Id}}}
\newcommand{\im}{\textrm{\textnormal{im}}}
\newcommand{\norm}[1]{\abs{\abs{#1}}}
\newcommand{\tpose}{^{T}\!}
\newcommand{\iprod}[1]{\left<#1\right>}
\newcommand{\giprod}{\iprod{\;\,,\;}}
\DeclareMathOperator{\tr}{tr}
\DeclareMathOperator{\trace}{tr}
\newcommand{\chgBasMat}[3]{\!\!\tensor*[_{#1}]{\left[#2\right]}{_{#3}}}
\newcommand{\vecBas}[2]{\tensor*[]{\left[#1\right]}{_{#2}}}
\DeclareMathOperator{\GL}{GL}
\DeclareMathOperator{\Mat}{Mat}
\DeclareMathOperator{\vspan}{span}
\DeclareMathOperator{\rank}{rank}
\newcommand{\V}[1]{\vec{#1}}
\DeclareMathOperator{\proj}{proj}
\DeclareMathOperator{\compProj}{comp}
\DeclareMathOperator{\row}{row}
\newcommand{\smallPMatrix}[1]{\paren{\begin{smallmatrix}#1\end{smallmatrix}}}
\newcommand{\smallBMatrix}[1]{\brack{\begin{smallmatrix}#1\end{smallmatrix}}}
\newcommand{\pmat}[1]{\begin{pmatrix}#1\end{pmatrix}}
\newcommand{\bmat}[1]{\begin{bmatrix}#1\end{bmatrix}}
\newcommand{\dual}{^{*}}
\newcommand{\pinv}{^{\dagger}}
\newcommand{\horizontalMatrixLine}{\ptxt{\rotatebox[origin=c]{-90}{$|$}}}
\DeclareMathOperator{\range}{range}
\DeclareMathOperator{\Symm}{Symm}
\DeclareMathOperator{\SU}{SU}

% Multilinear Algebra
\let\Lsym\L
\let\L\relax
\DeclareMathOperator{\L}{\mathscr{L}}
\DeclareMathOperator{\A}{\mathcal{A}}
\DeclareMathOperator{\Alt}{Alt}
\DeclareMathOperator{\Sym}{Sym}
\newcommand{\ot}{\otimes}
\newcommand{\ox}{\otimes}
\DeclareMathOperator{\asc}{asc}
\DeclareMathOperator{\asSet}{set}
\DeclareMathOperator{\sort}{sort}
\DeclareMathOperator{\ringA}{\mathring{A}}
\DeclareMathOperator{\Sh}{Sh}

% Topology
\newcommand{\closure}[1]{\overline{#1}}
\newcommand{\uball}{\mathcal{U}}
\DeclareMathOperator{\Int}{Int}
\DeclareMathOperator{\Ext}{Ext}
\DeclareMathOperator{\Bd}{Bd}
\DeclareMathOperator{\rInt}{rInt}
\DeclareMathOperator{\ch}{ch}
\DeclareMathOperator{\ah}{ah}
\newcommand{\LargerTau}{\mathlarger{\mathlarger{\mathlarger{\mathlarger{\tau}}}}}
\newcommand{\Tau}{\mathcal{T}}

% Analysis
\DeclareMathOperator{\Graph}{Graph}
\DeclareMathOperator{\epi}{epi}
\DeclareMathOperator{\hypo}{hypo}
\DeclareMathOperator{\supp}{supp}
\newcommand{\lint}[2]{\underset{#1}{\overset{#2}{{\color{black}\underline{{\color{white}\overline{{\color{black}\int}}\color{black}}}}}}}
\newcommand{\uint}[2]{\underset{#1}{\overset{#2}{{\color{white}\underline{{\color{black}\overline{{\color{black}\int}}\color{black}}}}}}}
\newcommand{\alignint}[2]{\underset{#1}{\overset{#2}{{\color{white}\underline{{\color{white}\overline{{\color{black}\int}}\color{black}}}}}}}
\newcommand{\extint}{\ptxt{ext}\int}
\newcommand{\extalignint}[2]{\ptxt{ext}\alignint{#1}{#2}}
\newcommand{\conv}{\ast}
\newcommand{\pd}[2]{\frac{\partial{}#1}{\partial{}#2}}
\newcommand{\del}{\nabla}
\DeclareMathOperator{\grad}{grad}
\DeclareMathOperator{\curl}{curl}
\let\div\relax
\DeclareMathOperator{\div}{div}
\DeclareMathOperator{\vol}{vol}

% Complex Analysis
\let\Re\relax
\DeclareMathOperator{\Re}{Re}
\let\Im\relax
\DeclareMathOperator{\Im}{Im}
\DeclareMathOperator{\Res}{Res}

% Abstract Algebra
\DeclareMathOperator{\ord}{ord}
\newcommand{\generated}[1]{\left<#1\right>}
\newcommand{\cycle}[1]{\smallPMatrix{#1}}
\newcommand{\id}{\ptxt{id}}
\newcommand{\iso}{\cong}
\DeclareMathOperator{\Aut}{Aut}
\DeclareMathOperator{\SL}{SL}
\DeclareMathOperator{\op}{op}
\newcommand{\isom}[4]{\!\!\!\begin{array}[t]{rcl}#1 & \!\!\!\!\overset{\sim}{\to} & \!\!\!\!#2\\ #3 & \!\!\!\!\mapsto & \!\!\!\!#4\end{array}}
\newcommand{\F}{\mathbb{F}}
\newcommand{\acts}{\;\reflectbox{\rotatebox[origin=c]{-90}{$\circlearrowright$}}\;}
\newcommand{\disjunion}{\mathrel{\text{\scalebox{1.07}{$\perp\mkern-10mu\perp$}}}}
\DeclareMathOperator{\SO}{SO}
\DeclareMathOperator{\stab}{stab}
\DeclareMathOperator{\nullity}{nullity}
\DeclareMathOperator{\Perm}{Perm}
\DeclareMathOperator{\nsubgp}{\vartriangleleft}
\DeclareMathOperator{\notnsubgp}{\ntriangleleft}
\newcommand{\presentation}[2]{\left<#1\;\middle|\;#2\right>}
\DeclareMathOperator{\Char}{char}
\DeclareMathOperator{\fchar}{char}
\DeclareMathOperator{\triv}{triv}
\DeclareMathOperator{\reg}{reg}
\DeclareMathOperator{\std}{std}
\DeclareMathOperator{\Func}{Func}
\DeclareMathOperator{\End}{End}

% Convex Optimization
\let\sectionSymbol\S
\let\S\relax
\newcommand{\S}{\mathbb{S}}
\DeclareMathOperator{\dist}{dist}
\DeclareMathOperator{\dom}{dom}
\DeclareMathOperator{\diag}{diag}
\DeclareMathOperator{\ones}{\mathbbm{1}}
\newcommand{\minimizeOver}[3]{\begin{array}{rl}\underset{#1}{\ptxt{minimize}} & #2\\ \ptxt{subject to} & #3\end{array}}
\newcommand{\maximizeOver}[3]{\begin{array}{rl}\underset{#1}{\ptxt{maximize}} & #2\\ \ptxt{subject to} & #3\end{array}}
\newcommand{\minimizationProblem}[2]{\minimizeOver{}{#1}{#2}}
\newcommand{\maximizationProblem}[2]{\maximizeOver{}{#1}{#2}}
\newcommand{\minimizeOverUnconstrained}[2]{\begin{array}{rl}\underset{#1}{\ptxt{minimize}} & #2\end{array}}
\newcommand{\maximizeOverUnconstrained}[2]{\begin{array}{rl}\underset{#1}{\ptxt{maximize}} & #2\end{array}}
\newcommand{\minimizationUnconstrained}[1]{\minimizeOverUnconstrained{}{#1}}
\newcommand{\maximizationUnconstrained}[1]{\maximizeOverUnconstrained{}{#1}}
\DeclareMathOperator{\argmin}{argmin}
\DeclareMathOperator{\argmax}{argmax}

% Proofs
\newcommand{\st}{s.t.}
\newcommand{\unique}{!}
\newcommand{\iffdef}{\overset{\ptxt{def}}{\Leftrightarrow}}
\newcommand{\eqVertical}{\rotatebox[origin=c]{90}{=}}
\newcommand{\mapsfrom}{\mathrel{\reflectbox{\ensuremath{\mapsto}}}}
\newcommand{\mapsdown}{\rotatebox[origin=c]{-90}{$\mapsto$}\mkern2mu}
\newcommand{\mapsup}{\rotatebox[origin=c]{90}{$\mapsto$}\mkern2mu}
\newcommand{\from}{\!\mathrel{\reflectbox{\ensuremath{\to}}}}
\newcommand{\labeledeq}[1]{\overset{\mathclap{\ptxt{#1}}}{=}}
\newcommand{\eqdef}{\labeledeq{def}}

% Brackets
\newcommand{\paren}[1]{\left(#1\right)}
\renewcommand{\brack}[1]{\left[#1\right]}
\renewcommand{\brace}[1]{\left\{#1\right\}}
\newcommand{\ang}[1]{\left<#1\right>}

% Algorithms
\algrenewcommand{\algorithmiccomment}[1]{\hskip 1em \texttt{// #1}}
\algrenewcommand\algorithmicrequire{\textbf{Input:}}
\algrenewcommand\algorithmicensure{\textbf{Output:}}
\newcommand{\algP}{\ptxt{\textbf{P}}}
\newcommand{\algNP}{\ptxt{\textbf{NP}}}
\newcommand{\algNPC}{\ptxt{\textbf{NP-Complete}}}
\newcommand{\algNPH}{\ptxt{\textbf{NP-Hard}}}
\newcommand{\algEXP}{\ptxt{\textbf{EXP}}}
\DeclareMathOperator{\fl}{fl}

%%%%%%%%%%%%%%%%%%%%%%%%%%%%%
% Other commands
%%%%%%%%%%%%%%%%%%%%%%%%%%%%%
\newcommand{\flag}[1]{\textbf{\textcolor{red}{#1}}}
\let\uSym\u
\let\u\relax
\newcommand{\u}[1]{\underline{#1}}
\let\bSym\b
\let\b\relax
\newcommand{\b}[1]{\textbf{#1}}
\let\iSym\i
\let\i\relax
\newcommand{\i}[1]{\textit{#1}}
\let\scSym\sc
\let\sc\relax
\newcommand{\sc}[1]{\textsc{#1}}

%%%%%%%%%%%%%%%%%%%%%%%%%%%%%%%%%%%%%%%
% Make l's curvy in math environments %
%%%%%%%%%%%%%%%%%%%%%%%%%%%%%%%%%%%%%%%
\mathcode`l="8000
\begingroup
\makeatletter
\lccode`\~=`\l
\DeclareMathSymbol{\lsb@l}{\mathalpha}{letters}{`l}
\lowercase{\gdef~{\ifnum\the\mathgroup=\m@ne \ell \else \lsb@l \fi}}%
\endgroup

%%%%%%%%%%%%%%%%%%%%%%%%%
% Fix \vdots and \ddots %
%%%%%%%%%%%%%%%%%%%%%%%%%
\usepackage{letltxmacro}
\LetLtxMacro\orgvdots\vdots
\LetLtxMacro\orgddots\ddots

\makeatletter
\DeclareRobustCommand\vdots{%
	\mathpalette\@vdots{}%
}
\newcommand*{\@vdots}[2]{%
	% #1: math style
	% #2: unused
	\sbox0{$#1\cdotp\cdotp\cdotp\m@th$}%
	\sbox2{$#1.\m@th$}%
	\vbox{%
		\dimen@=\wd0 %
		\advance\dimen@ -3\ht2 %
		\kern.5\dimen@
		% remove side bearings
		\dimen@=\wd2 %
		\advance\dimen@ -\ht2 %
		\dimen2=\wd0 %
		\advance\dimen2 -\dimen@
		\vbox to \dimen2{%
			\offinterlineskip
			\copy2 \vfill\copy2 \vfill\copy2 %
		}%
	}%
}
\DeclareRobustCommand\ddots{%
	\mathinner{%
		\mathpalette\@ddots{}%
		\mkern\thinmuskip
	}%
}
\newcommand*{\@ddots}[2]{%
	% #1: math style
	% #2: unused
	\sbox0{$#1\cdotp\cdotp\cdotp\m@th$}%
	\sbox2{$#1.\m@th$}%
	\vbox{%
		\dimen@=\wd0 %
		\advance\dimen@ -3\ht2 %
		\kern.5\dimen@
		% remove side bearings
		\dimen@=\wd2 %
		\advance\dimen@ -\ht2 %
		\dimen2=\wd0 %
		\advance\dimen2 -\dimen@
		\vbox to \dimen2{%
			\offinterlineskip
			\hbox{$#1\mathpunct{.}\m@th$}%
			\vfill
			\hbox{$#1\mathpunct{\kern\wd2}\mathpunct{.}\m@th$}%
			\vfill
			\hbox{$#1\mathpunct{\kern\wd2}\mathpunct{\kern\wd2}\mathpunct{.}\m@th$}%
		}%
	}%
}
\makeatother

\newcommand{\B}{
	\begin{tikzpicture}
	\filldraw [fill=red, draw=black] (0, 0) rectangle (0.37, 0.45);
	\draw [line width=0.5mm, white ] (0.1,0.08) -- (0.1,0.38);
	\draw[line width=0.5mm, white ] (0.1, 0.35) .. controls (0.2, 0.35) and (0.4, 0.2625) .. (0.1, 0.225);
	\draw[line width=0.5mm, white ] (0.1, 0.225) .. controls (0.2, 0.225) and (0.4, 0.1625) .. (0.1, 0.1);
	\end{tikzpicture}
}

% Allow custom symbols as arrows in tikz-cd
\tikzset{
  symbol/.style={
    draw=none,
    every to/.append style={
      edge node={node [sloped, allow upside down, auto=false]{$#1$}}}
  }
}
% Example: \arrow[r,symbol=\cong]
% https://tex.stackexchange.com/questions/394154/how-to-include-inclusion-subgroup-relationship-in-tikz-cd-diagram

\author{Thomas Cohn}
\title{Math 591 Lecture 34}
\date{11/20/20} % Can also use \today

\begin{document}
\maketitle
\setlength\RaggedRightParindent{\parindent}
\RaggedRight

\section*{Lie Derivatives of Forms}

\ex{
	Let $X=\ang{F^{1},\ldots,F^{n}}\in\mathfrak{X}(\R^{n})$, $\mu=dx^{1}\wedge\cdots\wedge{}dx^{n}$, and $\phi$ the flow of $X$. Then
	\[
		\mathcal{L}_{X}(\mu)=\restr{\frac{d}{dt}\phi_{t}^{*}\mu}{t=0}
	\]
	Well, $\phi_{t}^{*}\mu=\det(J(\phi_{t})_{*})\mu$, so
	\[
		\restr{\det(J(\phi_{t})_{*})}{t=0}=\tr\paren{\restr{\vphantom{\frac{d}{dt}}\smash{\underbrace{\frac{d}{dt\vphantom{a_{b_{c}}}}J(\phi_{t})_{*}}_{\star}}}{t=0}}\underbrace{\det(\phi_{t=0})}_{=1}
	\]
	For $\star$, do $\frac{d}{dt}$ first, and then $\pd{}{x^{i}}$. And $\det(\phi_{t=0})=1$, because $\phi_{t}$ is the flow of $X$, so $\frac{d}{dt}$ is just $X$. Thus, we have
	\[
		\tr \begin{pmatrix}
			\horizontalMatrixLine & \del{}F_{1} & \horizontalMatrixLine\\
			& \vdots\\
			\horizontalMatrixLine & \del{}F_{n} & \horizontalMatrixLine
		\end{pmatrix}=\sum_{j=1}^{n}\pd{F_{j}}{x^{j}}=\div(X)
	\]
	We conclude that $\mathcal{L}_{X}\mu=(\div{}X)\mu$.\proven
}

\thm{
	(Cartan's Magic Formula) $\mathcal{L}_{X}=\iota_{X}\of{}d+d\of\iota_{X}:\Omega^{k}\to\Omega^{k}$.\nn
	Proof: Let $P_{X}=\iota_{X}\of{}d-d\of\iota_{X}$. Then $P_{X}$ has the following properties:
	\begin{enumerate}[leftmargin=4\parindent, topsep=0pt, itemsep=0pt, label=\arabic*)]
		\item $\R$-linearity
		\item It's a derivation w.r.t. $\wedge$: $P_{X}(\alpha\wedge\beta)=P_{X}(\alpha)\wedge\beta+\alpha\wedge{}P_{X}(\beta)$
		\item It commutes with $d$
		\item $f\in{}C^{\infty}(M)\Rightarrow{}P_{X}(f)=X(f)$
		\item $P_{X}$ is local
	\end{enumerate}\up\n
	These properties belong to a unique operator. (Compute in coordinates, and by linearity, just use monomials.)
	\[
		P_{X}(a\,dx^{I})=P_{X}(a)\wedge{}dx^{I}+a\wedge{}P_{X}(dx^{I})=X(a)\,dx^{I}+aP_{X}(dx^{I})
	\]
	We then expand $P_{X}(dx^{I})$ with induction.\proven
}

\ex{
	For $k=2$:
	\[
		P_{X}(dx^{1}\wedge{}dx^{2})=P_{X}(dx^{1})\wedge{}dx^{2}+dx^{1}\wedge{}P(dx^{2})=dX(x^{1})\wedge{}dx^{2}+dx^{1}\wedge{}dX(x^{2})
	\]
	This is unique, and it just uses the five properties.\n
}

\subsection*{Applications of Cartan's Formula}

\defn{
	A \u{symplectic manifold} is a pair $(M,\omega)$, with $M$ a manifold, and $\omega\in\Omega^{2}(M)$ such that
	\begin{enumerate}[leftmargin=4\parindent, topsep=0pt, itemsep=0pt, label=\arabic*)]
		\item $\forall{}p\in{}M,v\in{}T_{p}M\cut\set{0}$, $\iota_{v}(\omega_{p})=\omega_{p}(v,\cdot):T_{p}M\to\R$ is nonzero.
		\item $d\omega=0$.
	\end{enumerate}\up\n
}

\par\noindent
Question: What are the symmetries of a symplectic manifold $(M,\omega)$? Specifically, are there one-parameter groups $\varphi_{t}:M\to{}M$ such that $\forall{}t$, $\varphi_{t}^{*}\omega=\omega$?\n
Use the Lie derivative: If $X\in\mathfrak{X}(M)$ is the generator of $\varphi_{t}$, then
\[
	\varphi_{t}^{*}\omega=\omega\Leftrightarrow\mathcal{L}_{X}\omega=0\Leftrightarrow\iota_{X}\underbrace{d\omega}_{=0}+d\iota_{X}\omega=0\Leftrightarrow{}d\iota_{X}\omega=0
\]

\defn{
	One particular class of such $X$'s comes from the following: Take any function $H\in{}C^{\infty}(M)$, and define $X$ by $\iota_{X}\omega=-dH$. By non-degeneracy of $\omega$, $X$ is unique! And,
	\[
		\map{T_{p}M}{T_{p}^{*}M}{v}{\iota_{v}\omega}
	\]
	has no kernel, so it's a bijection. This $X$ is called the \u{Hamilton field} of $H$.\n
}

\exer{
	Take $M=\R^{2n}$ with coordinates $(x,p)$, where $x$ and $p$ are the standard coordinates in $\R^{n}$. Let
	\[
		\omega=\sum_{i=1}^{n}dp^{i}\wedge{}dx^{i}\qquad{}H=\frac{1}{2}\norm{p}^{2}+V(x)
	\]
	Compute the Hamilton field, and show that the integral curves satisfy $\ddot{x}(t)=-\del{}V(x(t))$. This is better known as Newton's second law!\n
}

\par\noindent
We're now done with Lie derivatives.

\section*{Integration of Forms on Oriented Manifolds}

\par\noindent
First, we have to define orientation. Let $V$ be an $n$-dimensional vector space. Let $\mathscr{B}(V)$ be the set of all ordered bases of $V$. For $e\in\mathscr{B}$, $e=(e_{1},\ldots,e_{n})$ is an ordered basis of $V$. Observe: $\forall{}e,f\in\mathscr{B}$, $\exists\unique{}M\in\GL(n,\R)$ such that $\forall{}i$, $e_{i}=Mf_{i}$.\n

\defn{
	$e\sim{}f\Leftrightarrow\det{}M>0$. We say $e$ and $f$ define the same orientation of $V$.\n
}

\par\noindent
Check:
\begin{enumerate}[topsep=0pt, itemsep=0pt, label=\arabic*)]
	\item $\sim$ is an equivalence relation
	\item $\mathscr{B}/\sim$ has two elements
\end{enumerate}\up\n

\defn{
	An \u{orientation} of $V$ is a choice of an equivalence class in $\mathscr{B}/\sim$. Bases in that equivalence class are said to be \u{positive}.\n
}

\par\noindent
Alternatively, consider the set of nonzero top-degree forms, $\underbrace{\paren{{\bigwedge}^{n}V}}_{\dim=1}\cut\set{0}$. When we take away $0$, any $\mu\in\paren{{\bigwedge}^{n}V}\cut\set{0}$ defines an orientation by: a basis $e$ is positive iff $\mu(e)>0$. The orientation defined by $\mu$ only depends on which connected component of $\paren{\bigwedge^{n}V}\cut\set{0}$ contains $\mu$. Conversely, an orientation singles out one of the two components of $\paren{\bigwedge^{n}V}\cut\set{0}$.\n

\par\noindent
Conclusion: An orientation is a choice of a connected component of $\paren{\bigwedge^{n}V}\cut\set{0}$.\n

\par\noindent
Now, we move on to manifolds! Note: Not all manifolds are orientable (e.g. the Mobius band).\n

\defn{
	An \u{orientation} on $M$ (if it exists) is a continuous choice of orientation of each tangent space.\n
}

\par\noindent
Continuity means $\forall{}p\in{}M$, there exists a continuous moving frame $(E_{1},\ldots,E_{n})$ such that at every point $q$ in the domain of $E$, $(E_{1}(q),\ldots,E_{n}(q))$ is a positive basis of $T_{q}M$.\n

\lemma{
	A connected manifold can have either two orientations, or it's non-orientable.\nn
	Proof: The idea is if the manifold is orientable, then consider orientations $\mathscr{O}_{1}$ and $\mathscr{O}_{2}$. Define $F:M\to\R$ by $f(p)=1$ if $\mathscr{O}_{1}(p)=\mathscr{O}_{2}(p)$, and $0$ otherwise. (Note: we don't really ever use this notation.) Then by the continuity of $\mathscr{O}_{1}$ and $\mathscr{O}_{2}$, $f$ is continuous, so $f$ is locally constant, so if $M$ is connected, then $f$ is constant.\proven
}

\defn{
	Let $M$ be an oriented manifold. Then a \u{positive atlas} on $M$ is an atlas $\set{(U_{\alpha},\phi_{\alpha})}$ of $M$ such that $\forall\alpha$, the moving frame $\paren{\pd{}{x^{1}_{\alpha}},\ldots,\pd{}{x^{n}_{\alpha}}}$ is a positive frame.\n
}

\newpage
\lemma{
	\begin{enumerate}[leftmargin=4\parindent, topsep=0pt, itemsep=0pt, label=\arabic*)]
		\item The transition functions $F$ and $G$ of any two elements in a positive atlas satisfy $\det(J(F))=1=\det(J(G))$.
		\item An oriented manifold always has a positive atlas.
	\end{enumerate}\up\nn
	Proof: This is very tedious! Idea:\n
	1) Recall that $J(F)$ is actually the change of basis matrix between $\paren{\pd{}{x^{1}_{\alpha}},\ldots,\pd{}{x^{n}_{\alpha}}}$.
	2) Start with any atlas. $\forall\alpha$, $\set{\pd{}{x^{1}_{\alpha}},\ldots,\pd{}{x^{n}_{\alpha}}}$ is either positive, or not. If it is positive, do nothing, and keep $\phi_{\alpha}$. If it's not positive, relable (switch) $x^{1}$ and $x^{2}$. Now it's positive!\n
	\proven
}

\end{document}